% ===== 底本 第 5 章《环簇的 SYZ 镜像对称（续）》；片段 tr06a；源块 B189--B214 =====

\section{环簇的 SYZ 镜像对称（续）}

% [源 B189]
本讲给出 [CO06] 的更多细节。让我们从环簇流形的回顾开始。

% [Tutor Note: B189 原文作 "recollment"，当为 "recollection"（回顾）之误；据上下文译作「回顾」。]

% [源 B190]
\subsection{环簇流形：复结构与辛结构}

% [源 B191]
设 $N$ 为格（lattice），$M$ 为其对偶格（dual lattice），并记 $\langle -,-\rangle\colon M\times N\to\mathbb{Z}$ 为赋值配对（evaluation pairing）。

% [源 B192]
\begin{definition*}
一个 $k$ 维的\emph{正则锥}（regular cone）是 $\sigma\subseteq N_{\mathbb{R}}=N\otimes_{\mathbb{Z}}\mathbb{R}$，其生成元满足 $\langle v_1,\dotsb, v_k\rangle_{\geq\mathbb{R}}=\sigma$。$\sigma$ 的\emph{面}（face）是 $\langle v_{i_1},\dotsb, v_{i_r}\rangle_{\geq\mathbb{R}}$。$1$ 维锥称为\emph{射线}（ray）。
\end{definition*}

% [源 B193]
\begin{definition*}
一个\emph{扇}（fan）是 $\Sigma=\{\text{regular cones}\}$，使得
\begin{enumerate}[label=\arabic*),leftmargin=1.9em,itemsep=1pt,topsep=2pt]
  % [源 B194]
  \item 对任意 $\sigma\in\Sigma$，$\sigma$ 的任意面都落在 $\Sigma$ 中；
  % [源 B195]
  \item 对给定的 $\sigma,\sigma'\in\Sigma$，交 $\sigma\cap\sigma'$ 同时是 $\sigma$ 与 $\sigma'$ 的面。
\end{enumerate}
\end{definition*}

% [源 B196]
称 $\Sigma$ 为\emph{完备的}（complete），若 $N_{\mathbb{R}}=\displaystyle\bigcup_{\sigma\in\Sigma}\sigma$。记 $\Sigma^{(k)}=\{\text{all }k\text{-dimensional cones in }\Sigma\}$，以及 $G(\Sigma)=\{\text{generators of rays}\}$。

% [源 B197]
\begin{definition*}
子集 $P\subseteq G(\Sigma)$ 称为\emph{原语集}（primitive collection），若 $P$ 本身不构成锥，而 $P$ 的任意真子集都构成锥。记 $P(\Sigma)=\{\text{all primitive collections}\}$。
\end{definition*}

% [源 B198]
\begin{example*}
设 $G(\Sigma)=\{e_1,\dotsb,e_n,-e_1-\dotsb-e_n\}$。任何不等于 $G(\Sigma)$ 的集合都生成一个锥，故 $G(\Sigma)$ 是原语集。
\end{example*}

% [源 B199]
考虑 $\mathbb{C}^N$，其坐标 $z_1,\dotsb, z_N$ 对应 $G(\Sigma)$ 中的向量 $v_1,\dotsb, v_N$。

% [源 B200]
\begin{definition*}
对一个原语集 $P$，记 $\mathbb{A}(P)=\{(z_1,\dotsb, z_N)\colon z_i=0\text{ if }v_i\in P\}$。
\end{definition*}

% [Tutor Note: B193、B196、B197、B200、B204 公式内 \text{...} 的英文说明系公式源码的一部分（如「regular cones」「all primitive collections」「if」「for」），按「公式逐字转写」的要求保留英文未译；正文部分均译为中文。]

% [源 B201]
在上一个例子中，$\mathbb{A} (G(\Sigma))=\{(0,0,\dotsb, 0)\}$，且商 $\left(\mathbb{C}^{n+1}\setminus\mathbb{A} (G(\Sigma))\right)\big/\mathbb{C}^{\times}\cong\mathbb{P}^n$。记 $Z(\Sigma)=\displaystyle\bigcup_{P\in P(\Sigma)}\mathbb{A}(P)$，$U(\Sigma)=\mathbb{C}^N\setminus Z(\Sigma)$。映射 $e_i\mapsto v_i$ 给出群同态 $\pi\colon\mathbb{Z}^N\to\mathbb{Z}^n$，其核为 $\mathbb{K}=\ker (\pi)$。每个向量 $\vec{\lambda} =(\lambda_1,\dotsb,\lambda_N)\in\mathbb{K}$ 通过 $a_{\vec{\lambda}}(t).(z_1,\dotsb, z_N)=(t^{\lambda_1}z_1,\dotsb, t^{\lambda_N}z_N)$ 作用在 $\mathbb{C}^N$ 上。令 $D(\Sigma)=\langle a_{\vec{\lambda}}\colon\vec{\lambda}\in\mathbb{K}\rangle\subseteq\left(\mathbb{C}^{\times}\right)^N$ 为 $\mathbb{K}$ 生成的代数环面（algebraic torus）。

% [源 B202]
\begin{lemma*}
$D(\Sigma)$ 自由地作用在 $U(\Sigma)$ 上。
\end{lemma*}

% [源 B203]
与扇 $\Sigma$ 相伴的\emph{环簇}（toric variety）定义为 $X_{\Sigma}=U(\Sigma)/D(\Sigma)$。在上一个例子中，映射 $\pi$ 由 $\mathbb{Z}^{n+1}\to\mathbb{Z}^n\to 0$ 给出，其中 $e_i\mapsto e_i$（$1\leq i\leq n$）、$e_{n+1}\mapsto -e_1-\dotsb -e_n$，于是 $\mathbb{K} =\langle e_1+\dotsb +e_n\rangle$，并且所得的作用 $a_{\lambda}(t).(z_1,\dotsb, z_n)=(t^{\lambda}z_1,\dotsb, t^{\lambda}z_n)$ 就是 $\mathbb{P}^n$ 上通常的对角作用。

% [Tutor Note: B203 中 $\mathbb{K} =\langle e_1+\dotsb +e_n\rangle$ 疑误：按 $\pi\colon\mathbb{Z}^{n+1}\to\mathbb{Z}^n$ 的上述定义，核应由 $e_1+\dotsb +e_n+e_{n+1}$ 生成；同一句中作用式作用于 $\mathbb{C}^{n+1}$，其坐标写成 $(z_1,\dotsb, z_n)$ 也应作 $(z_1,\dotsb, z_{n+1})$。均照原文转写，未改。]

% [源 B204]
\begin{definition*}
对正则锥 $\sigma=\langle v_{i_1},\dotsb,v_{i_r}\rangle$，记 $U_{\sigma}=\{(z_1,\dotsb,z_N)\in\mathbb{C}^N\vert z_j\neq 0\text{ for }j\not\in\{i_1,\dotsb, i_r\}\}$。
\end{definition*}

% [源 B205]
对 $\mathbb{P}^n$，我们得到 $U_i=\{z_i\neq 0\}$。

% [源 B206]
\begin{proposition*}
(1) $U(\Sigma)=\displaystyle\bigcup_{\sigma} U(\sigma)$；(2) 若 $\sigma$ 是 $\sigma'$ 的面，则 $U(\sigma)\supseteq U(\sigma')$；(3) $U(\sigma_1\cap\sigma_2)=U(\sigma_1)\cap U(\sigma_2)$。
\end{proposition*}

% [Tutor Note: B206 的 (2) 方向疑误：按 B204 中 $U_{\sigma}$ 的定义，当 $\sigma$ 是 $\sigma'$ 的面时（$I_\sigma\subset I_{\sigma'}$，条件更少）应有 $U(\sigma)\subseteq U(\sigma')$，而原文作 $\supseteq$；以 $\mathbb{P}^1$ 的扇 $G=\{e_1,-e_1\}$ 取 $\sigma=0$ 可验证。照译原文，已登记 README「原文笔误登记表」。]

% [源 B207]
\begin{proposition*}
设 $\sigma=\langle v_1,\dotsb, v_n\rangle$，取其生成元的对偶基 $\{u_{i_1},\dotsb, u_{i_n}\}$，则 $U(\sigma)\cong\mathbb{C}^n\times\left(\mathbb{C}^{\times}\right)^{N-n}$，其中 $U_{\sigma}=U(\sigma)/D(\Sigma)\cong\mathbb{C}^n$，其坐标由下式给出：
\begin{gather*}
  x_1^{\sigma}=z_1^{\langle v,u_{i_1}\rangle}\dotsb z_N^{\langle v_N,u_{i_1}\rangle}\\
  \vdots\\
  x_n^{\sigma}=z_1^{\langle v_1,u_{i_n}\rangle}\dotsb z_N^{\langle v_N,u_{i_n}\rangle}
\end{gather*}
\end{proposition*}

% [Tutor Note: B207 第一行 $z_1^{\langle v,u_{i_1}\rangle}$ 中的 $\langle v,u_{i_1}\rangle$ 似漏写下标（对照第三行 $\langle v_1,u_{i_n}\rangle$，此处当为 $\langle v_1,u_{i_1}\rangle$）；又「对偶基 $\{u_{i_1},\dotsb, u_{i_n}\}$」的下标记号与 $v_1,\dotsb,v_n$ 不对应，疑本应作 $\{u_1,\dotsb,u_n\}$。均照原文转写，未改。]

% [源 B208]
\subsection{辛几何方面}

% [Tutor Note: B208 原文小节标题作 "Sympletic aspects."，其中 Sympletic 为 Symplectic（辛的）之拼写错误；按 BRIEF §3 定稿译名作「辛几何方面」。]

% [源 B209]
考虑多面体 $P=\{x\in M_{\mathbb{R}}\colon\langle x,v_i\rangle\geq\lambda_i\}$。对一个面 $\Gamma$，记 $\sigma_{\Gamma}\colon =\displaystyle\bigcup_{r\geq 0} r\cdot (P -m)$，其中 $m\in\mathrm{int}\ \Gamma$。于是得到\emph{对偶锥}（dual cone）$\check{\sigma}_{\Gamma} =\{x\in N_{\mathbb{R}}\colon\langle y,x\rangle\geq 0\ \forall y\in\sigma_{\Gamma}\}$ 与诱导出的完备扇 $\Sigma (P)$。$\left(\mathbb{C}^{\times}\right)^N$ 的代数环面作用给出实环面 $T^N$ 在 $U(\Sigma (P))$ 上的辛作用，因此得到矩映射（moment map）$\mu\colon U(\Sigma (P))\to T_e^{\ast} T^N\ni r$，从而\emph{辛环簇流形}（symplectic toric manifold）就是辛约化（symplectic reduction）$\mu^{-1} (r)/T^N=X_{\Sigma(P)}$，连同商辛结构。

% [Tutor Note: B209 中矩映射取值于 $T_e^{\ast} T^N\cong\mathbb{R}^N$（李代数的对偶），$r$ 即与多面体 $P$ 对应的正则值；$\mu^{-1}(r)/T^N$ 是 Marsden--Weinstein 辛约化，作者省略了这一步的论证。]

% [源 B210]
\subsection{Morse--Bott Floer 上同调}

% [源 B211]
本节讨论拉格朗日子流形的 Morse--Bott 版本 Floer 上同调。设 $M$ 为光滑流形。

% [Tutor Note: B211 原文 "smoothmanifold" 缺一空格，应为 "smooth manifold"（光滑流形），按后者理解。]

% [源 B212]
% [Tutor Note: B212 原文作 "Definition".（带引号），表示此处是宽松/非严格意义下的定义；本译本按 definition* 排印。]
\begin{definition*}
一个 $k$ 维的\emph{几何链}（geometric chain）是一个 $k$ 维单纯复形（simplicial complex）$P$，满足 $P=P_{reg}\cup P_S$，其中 $P_{reg}$ 是一个 $k$ 维有向子流形并满足若干条件，$f\colon P\to M$ 是连续函数，且 $f\vert_{P_{reg}}$ 光滑。
\end{definition*}

% [源 B213]
对几何链 $[P,f]$，我们得到函数
\begin{align*}
  T_{[P,f]}\colon\Omega_c^k (M)&\to \mathbb{R}\\
  \alpha&\mapsto \int_{P_{reg}} f^{\ast}\alpha
\end{align*}
定义微分 $\partial T_{[P,f]} (\alpha) =T_{[P,f]} (d\alpha)$。这给出上同调 $H_{Sing}^{\ast} (M)$。若 $[P,f]$ 满足 $\dim f(P)<\dim P$，则 $[P,f]$ 是 $\Omega^{\ast} (M)$ 上的零映射。

% [Tutor Note: B213 的 $\partial T_{[P,f]}(\alpha)=T_{[P,f]}(d\alpha)$ 即边缘算子（$\partial$ 把 $k$ 次流降到 $k-1$ 次；$T_{[P,f]}$ 为 $\Omega_c^k(M)$ 上的线性泛函）；「$\dim f(P)<\dim P$ 时为零映射」成立的理由是：此时 $f^{\ast}\alpha$ 作为 $P_{reg}$ 上的 $k$ 次形式必为零，故积分恒为 $0$。]

% [源 B214]
现在考虑配对 $(M,L)$，其中 $M$ 是辛流形，$L\subseteq M$ 是 $M$ 的拉格朗日子流形（Lagrangian submanifold）。对 $\beta\in\pi_2 (M,L)$，可以定义
\[
\mathrm{Aut}\ (D)\curvearrowright\tilde{\mathcal{M}}_n (L,J;\beta)=\left\{\begin{array}{l}
  w\colon (\mathbb{D},\partial\mathbb{D};z_1,\dotsb, z_n)\to (M,L)\\
  \bar{\partial}_Jw=0
\end{array}\right\}
\]
而商 $\mathcal{M}_n =\tilde{\mathcal{M}}_n/\mathrm{Aut}\ (D)$ 具有紧化 $\overline{\mathcal{M}}_n$。

% [Tutor Note: B214 中「具有紧化」的论证即 Gromov 紧致性（底本参考文献中的 [MS04]），作者此处未展开；公式里的 $z_1,\dotsb, z_n$ 是圆盘 $\mathbb{D}$ 上的 $n$ 个标记点。]
